\chapter{Derivations and Transport Limits}\label{app:derivations}

\section{Uniform two-way relaxation}
For the uniform system in \cref{eq:two_way_test}, define slip $w=v-u$.
Adding the fluid and population-weighted particle equations gives
\begin{equation}
 \frac{d}{dt}(u+\beta v)=0.
\end{equation}
Subtracting the equations gives
\begin{equation}
 \dot w=-\frac{1+\beta}{\St}w,\qquad
 w(t)=w_0\exp\left[-\frac{1+\beta}{\St}t\right].
\end{equation}
If $M_0=u_0+\beta v_0$ is the conserved momentum, the velocities are
\begin{equation}
 u(t)=\frac{M_0-\beta w(t)}{1+\beta},\qquad
 v(t)=\frac{M_0+w(t)}{1+\beta}.
\end{equation}
For $u_0=0$, $v_0=1$ and $\beta=0.25$, both velocities tend to $0.2$.
The total kinetic energy is
\begin{equation}
 K=\frac{1}{2}u^2+\frac{\beta}{2}v^2,
 \qquad \dot K=-\frac{\beta}{\St}(v-u)^2.
\end{equation}
This gives the dissipative reference behavior of V02. It concerns a
closed, uniform exchange system; the driven channel has additional work
and boundary terms.

\section{Adjoint transfer and kinetic energy}
Let $J$ denote the interpolation operator with entries
$J_{ia}=W_{ai}\Delta A_a$. With quadrature mass matrix
$M_A=\operatorname{diag}(\Delta A_a)$, the spreading operator is
$-M_A^{-1}J^{\mathsf T}$ when the parcel force vector stores integrated
forces. Thus
\begin{equation}
 \vect{U}=J\vect{u},\qquad
 \vect{S}=-M_A^{-1}J^{\mathsf T}\vect{F}.
\end{equation}
The fluid work is
$\vect{u}^{\mathsf T}M_A\vect{S}
=-(J\vect{u})^{\mathsf T}\vect{F}$.
Adding parcel drag work produces
\cref{eq:drag_dissipation}. This matrix statement makes explicit why
interpolation and spreading must use the same quadrature and weights.

For a collision-free driven channel with no-slip stationary walls and
periodic streamwise faces, a formally exact solution obeys
\begin{equation}
 \begin{aligned}
 \frac{d}{dt}\left[\frac{1}{2}\int_\Omega|\vect{u}|^2\,dA
 +\sum_i\frac{m_i}{2}|\vect{v}_i|^2\right]
 &=\int_\Omega\vect{u}\cdot\vect{f}_d\,dA
 +\sum_i m_i\vect{v}_i\cdot\vect{g}_{\mathrm{eff}}\\
 &\quad-\frac{1}{\Rey_H}\int_\Omega|\nabla\vect{u}|^2\,dA
 -\sum_i\frac{m_i}{\St_i}|\vect{U}_i-\vect{v}_i|^2.
 \end{aligned}
 \label{eq:energy_budget}
\end{equation}
Pressure and periodic-face contributions cancel under the stated
conditions. The wall velocity is zero, so stationary walls do no work.
Quadrature and neural equation residuals can create a nonzero discrepancy
when this identity is evaluated numerically. Wall tractions still enter
the momentum budget; they must not be omitted merely because wall work
vanishes in the energy budget.

\section{A limiting pressure-drop estimate}
For slowly varying plane walls and negligible inertia, the local
lubrication approximation gives the dimensional velocity
\begin{equation}
 u(x,y)=-\frac{1}{2\mu}\frac{dp}{dx}[h(x)^2-y^2].
\end{equation}
Integrating over $-h\leq y\leq h$ yields
\begin{equation}
 Q=-\frac{2h(x)^3}{3\mu}\frac{dp}{dx},\qquad
 \Delta p=\frac{3\mu Q}{2}\int_{x_a}^{x_b}\frac{dx}{h(x)^3}.
 \label{eq:lubrication_drop}
\end{equation}
For constant half-height $h_0$, this reduces to the plane Poiseuille
pressure drop $3\mu Q(x_b-x_a)/(2h_0^3)$. At fixed flux, the normalized
hydraulic resistance estimate is
\begin{equation}
 \frac{\Delta p}{\Delta p_0}
 =\frac{h_0^3}{x_b-x_a}\int_{x_a}^{x_b}\frac{dx}{h(x)^3}.
\end{equation}
This is a limiting analytical check, not a simulation result for the
proposed severity values. Its reliability requires sufficiently small wall
slope and inertial effects; a full Navier--Stokes solution is needed when
those assumptions fail \citep{batchelor1967}.

\section{Baseline trajectory criteria}\label{sec:trajectory_appendix}
\Cref{tab:path_limits} retains the seven baseline comparisons associated
with the field results in \cref{ch:basic}. RMS position discrepancy is
normalized by $W=0.3$. The current numerical limits require discrepancy
below 5\%, common-time coverage above 95\%, and agreement of every endpoint
class. Some individual rows satisfy those conditions, but the complete set
does not. The trajectories are consequently not used as validated
deposition evidence.

\begin{table}[htbp]\centering\small
\caption{Complete baseline particle comparison; no initial position is
removed to improve the aggregate outcome.}\label{tab:path_limits}
\begin{tabularx}{\textwidth}{@{}rrrY@{}}\toprule
Initial $\eta$ & RMS position/$W$ (\%) & Common time (\%) & Same endpoint class\\\midrule
\EvidenceRows{trajectories}
\bottomrule\end{tabularx}
\end{table}

\section{Planar mass convention}
For a physical sphere of diameter $d$, mass is $\rho_p\pi d^3/6$.
A parcel containing $N_i$ such spheres has mass
$N_i\rho_p\pi d^3/6$. Dividing by $\rho_fH^2B$ gives
\begin{equation}
 \frac{N_i\rho_p\pi d^3/6}{\rho_fH^2B}
 =\frac{N_iR\pi(d/H)^3}{6(B/H)},
\end{equation}
which is \cref{eq:parcel_mass}. The dimensional reaction source in a planar
fluid equation has force divided by depth and area. This is why a
two-dimensional Gaussian normalization uses area rather than volume. The
reference depth records the physical interpretation without introducing a
resolved spanwise flow field.
